\subsection{CP Symmetry and Symplectic Modular Invariance --- arXiv:2102.06716}

\textbf{Authors:} Gui-Jun Ding, Ferruccio Feruglio, Xiang-Gan Liu

\paragraph{主な主張}
シンプレクティック・モジュラー不変理論におけるCPの定義を分類し、種数 $g\geq3$ では一意、$g\leq2$ では独立な二通りが可能であることと、モジュライ空間内のCP保存領域を示した\cite{Ding:2021iqp}。

\paragraph{新規性}
多変数のSiegelモジュラー形式・物質場・モジュライに整合的に作用するCPを定式化し、少数のパラメータでレプトン混合とCP位相を記述する模型に適用した。

\paragraph{位置付け}
多モジュライのフレーバー模型でCPを課し、真空によるCP破れを判定するための基礎となる。
